% Part: many-valued-logic
% Chapter: sequent-calculus
% Section: structural-rules

\documentclass[../../../include/open-logic-section]{subfiles}

\begin{document}

\olfileid{mvl}{seq}{str}

\olsection{গঠনগত বিধি}

$n$-পাশবিশিষ্ট সিকোয়েন্ট কলনের গঠনগত বিধিগুলি ধ্রুপদি ক্ষেত্রের
মতোই কাজ করে, তবে প্রতিটি অবস্থান~$i$-এর জন্য পৃথকভাবে।

\begin{defish}
\begin{center}
\AxiomC{$ \Gamma_1 \nSequent \dots \nSequent \phantom{!A,}\Gamma_i \nSequent \dots \nSequent \Gamma_n $}
\RightLabel{$\iR{\Weakening}{i}$}
\UnaryInfC{$\Gamma_1 \nSequent \dots \nSequent !A, \Gamma_i \nSequent \dots \nSequent \Gamma_n$}
\DisplayProof
\\[2ex]
\AxiomC{$ \Gamma_1 \nSequent \dots \nSequent !A, !A, \Gamma_i \nSequent \dots \nSequent \Gamma_n $}
\RightLabel{$\iR{\Contraction}{i}$}
\UnaryInfC{$\Gamma_1 \nSequent \dots \nSequent \phantom{!A,}!A, \Gamma_i \nSequent
\dots \nSequent \Gamma_n$}
\DisplayProof
\\[2ex]
\AxiomC{$ \Gamma_1 \nSequent \dots \nSequent \Gamma_i, !A, !B, \Gamma_i' \nSequent \dots \nSequent \Gamma_n $}
\RightLabel{$\iR{\Exchange}{i}$}
\UnaryInfC{$\Gamma_1 \nSequent \dots \nSequent \Gamma_i, !B, !A, \Gamma_i' \nSequent \dots
\nSequent \Gamma_n$}
\DisplayProof
\end{center}
\end{defish}

দুর্বলীকরণ, সংকোচন ও অদলবদলের পরপর কয়েকটি অনুমানকে প্রায়ই
দ্বৈত অনুমান-রেখা দিয়ে নির্দেশ করা হবে।

\Cut{} বিধির কয়েকটি রূপ আছে: সিকোয়েন্টের পৃথক দুই অবস্থান
$i \neq j$ বেছে নেওয়ার প্রতিটি উপায়ের জন্য একটি রূপ:
\begin{defish}
\[
\AxiomC{$ \Gamma_1 \nSequent \dots \nSequent !A, \Gamma_i \nSequent \dots \nSequent \Gamma_n $}
\AxiomC{$ \Delta_1 \nSequent \dots \nSequent !A, \Delta_j \nSequent \dots \nSequent \Delta_n $}
\RightLabel{$\iR{\Cut}{i,j}$}
\BinaryInfC{$\Gamma_1,\Delta_1 \nSequent \dots \nSequent \Gamma_n, \Delta_n$}
\DisplayProof
\]
\end{defish}

\end{document}
